\documentclass{homework}
% \usepackage{graphicx} % Required for inserting images
\usepackage{amsmath}
\usepackage{gensymb}
\usepackage{subfig}
\usepackage{float}

\class{ECEN 5738}
\title{Homework 7}
\author{Sean Jennings}
\date{April 30, 2024}

\begin{document}
\maketitle

\section{Problem 1}

Define $\alpha: \R \to \R$ to be the common denominator of the dynamics equations.
That is, we rewrite the dynamics:
\begin{align*}
    \dot{x}_1 &= x_2     \\
    \dot{x}_2 &= \frac{1}{\alpha(x_3)} \biggl(
        -x_1 + \epsilon x_4^2 \sin x_3 - \epsilon \cos x_3 u
    \biggr)\\
    \dot{x}_3 &= x_4     \\
    \dot{x}_4 &= \frac{1}{\alpha(x_3)} \biggl(
        \epsilon \cos x_3 ( x_1 - \epsilon x_4^2 \sin x_3) + u
    \biggr) \\
\end{align*}
with
\begin{equation*}
    \alpha(x_3) = 1 - \epsilon^2 \cos^2 x_3
\end{equation*}

\begin{letterenum}
\item We differentiate the output $y = x_3$ to find the first derivative that depends
on $u$. Since $\dot{y} = \dot{x}_3 = x_4$, we have that the second
derivative $\ddot{y} = \dot{x}_4$ depends on $u$ and the relative degree is therefore
2.

\item We equate $\ddot{y}$ to a modified control variable $v\in\R$ and note that
$\dot{x}_4$ can be written in the form:
\begin{equation}
    \ddot{y} = \dot{x}_4 = v = \frac{1}{\alpha(x_3)} ( u + \phi(x))
    \label{eq:control-transformation}
\end{equation}
where $\phi: \R^4 \to \R$ is the functional:
\begin{equation*}
    \phi(x_1, x_2, x_3, x_4) = \epsilon \cos x_3 ( x_1 - \epsilon x_4^2 \sin x_3) 
\end{equation*}
Since $1 / \alpha > 0$, the control transformation (\ref{eq:control-transformation})
is invertible, $u$ can be expressed as a function of $v$
\begin{equation*}
    u = \alpha(x_3) v - \phi(x)
\end{equation*}
and the system is indeed input-output linearizable.

By letting $v = -\mat{y & \dot{y}} K$, we see the output dynamics are given by:
\begin{equation*}
    \mat{ \dot{y} \\ \ddot{y} } = \mat{\dot{y} \\ 0} + \mat{0 \\ 1} v =
        \mat{
        0 & 1 \\
        -k_1 & -k_2
    } \mat{y \\ \dot{y}}
\end{equation*}
which has $y=\dot{y} = 0$ as a globally asymptotically stable equilibrium point 
for appropriately chosen $K$.

\item The zero dynamics are found by taking $y = \dot{y} = \ddot{y} = 0$ which
implies that $x_3 = x_4 = v = 0$. Thus, $u = - \phi(x_1, x_2, 0, 0) = -\epsilon x_1$ and the dynamics
of the translational elements of oscillator are given by
\begin{align*}
    \dot{x}_1 &= x_2 \\
    \dot{x}_2 &= \frac{1}{\alpha(0)} \biggl(
        -x_1 - \epsilon u
    \biggr) \\
    &= \frac{1}{1 - \epsilon^2} (\epsilon^2 - 1) x_1 \\
    &= -x_1
\end{align*}
or in matrix form
\begin{equation*}
    \mat{\dot{x}_1 \\ \dot{x}_2} = \mat{
        0 & 1 \\
        -1 & 0
    } \mat{x_1 \\ x_2}
\end{equation*}

Therefore the zero dynamics are linear and marginally stable. Physically,
this corresponds to the pendulum being actively stabilized to its
downmost position while the cart oscillates due to the spring. Since the 
dynamics do not include damping/dissipative terms, this signifies that the control
is passive: it does not add energy to the system.

\end{letterenum}

\section{Problem 2}
Define $f$ as the polynomial:
\begin{equation*}
    f(\phi) = -\frac{3}{2} \phi^2 - \frac{1}{2} \phi^3
\end{equation*}
so that the dynamics of $\phi$ can be written as
\begin{equation*}
    \dot{\phi} = - \psi + f(\phi).
\end{equation*}
Consider $\psi$ as an input to the dynamics of $\phi$. Let $\alpha: \R \to \R$
be the function
\begin{equation*}
    \alpha(\phi) = k_1 \phi + f(\phi)
\end{equation*}
so that along the surface $(\phi, \psi) = (\phi, \alpha(\phi))$
we have $\dot{\phi} = -k_1 \phi$ which is stable for $k_1 > 0$.

Perform the change of variables $z = \psi - \alpha(\phi)$.
Then, the dynamics in the coordinates $(\phi, z)$ is given by:
\begin{align}
    \dot{\phi} &= - (z + \alpha(\phi)) + f(x) = -z - k_1 \phi \\
    \dot{z} &= \dot{\psi} - \frac{d\alpha}{d\phi} \dot{\phi}
    \label{eq:redefined-control}
\end{align}

We redefine the control variable $v$ to be equal to the right side of
(\ref{eq:redefined-control}):
\begin{align*}
v &= \dot{\psi} - \frac{d\alpha}{d\phi} \dot{\phi} \\
    &= \frac{1}{\beta}(\phi + 1 - u) - (k_1 - 3\phi - 3/2 \phi^2)
        ( - \psi - 3/2 \phi^2 - 1/2 \phi^3) \\
\end{align*}
We see that $\dot{z}$ is linear with respect to the control input $u$,
and can therefore define the function $g: \R^2 \to \R$ to be the terms which do
not depend on $u$, so that
\begin{equation*}
    v = -\frac{1}{\beta} u + g(\phi, \psi).
\end{equation*}

Finally, we use the Lyapunov function
\begin{equation*}
    V(\phi, z) = 1/2 (\phi^2 + z^2)
\end{equation*}
to find a control that stabilizes the origin. The 
derivative of $V$ is given by:
\begin{align*}
    \dot{V}(\phi, z) &= \phi \dot{\phi} + z \dot{z} \\
        &= \phi( - k_1 \phi - z) + zv \\ 
        &= -k_1 \phi^2 - (\phi - v) z
\end{align*}
so that choosing $v = \phi -k_2 z$ makes $\dot{V}$ globally
negative definite:
\begin{equation*}
    \dot{V}(\phi, z) = - k_1 \phi^2 - k_2 z^2,
\end{equation*}

Therefore, the control
\begin{align*}
    u &= \beta ( g(\phi, \psi) - v) \\
        &= \beta ( g(\phi, \psi) - \phi + k_2 z) \\
        &= \beta ( g(\phi, \psi) - \phi + k_2 (\psi - \alpha(\phi)))
\end{align*}
stabilizes the origin.

\section{Problem 3}
Define the surface
\begin{equation*}
    s(t) = ax_1 + x_2
\end{equation*}
On the surface $s(t) = 0$, we have
\begin{equation*}
    \dot{x}_1 = (1 - 3a) x_1
\end{equation*}

So for $a > 1/3$, $x_1$ converges to 0, and since the surface
crosses the $x_1$ axis through the origin any trajectory which stays on
the surface will approach the origin. Thus, a controller
which drives the system to $s$ in finite time will make the origin
asymptotically stable.

% Define the functional $f: \R^2 \to \R$ to be
% \begin{equation*}
%     f(x) = \phi x_2^4 + x_1
% \end{equation*}
% so that the $x_2$ dynamics can be written
% \begin{equation*}
%     \dot{\phi_2} = f(x) + u
% \end{equation*}
% Letting $\hat{f}$ be our best estimate of the dynamics:
% \begin{equation*}
%     \hat{f}(x) = \hat{\theta} x_2^4 + x_1.
% \end{equation*}
% where $\hat{\theta} = 1$. Then, the uncertainty in the dynamics
% $\tilde{f}(x) = f(x) - \hat{f}(x)$ can be bounded:
% \begin{equation*}
%     |tilde{f}(x)| \leq 0.2 x_2^4
% \end{equation*}

The surface dynamics are given by:
\begin{align*}
    \dot{s} &= a\dot{x}_1 + \dot{x}_2 \\
        &= a(x_1 + 3x_2) + \theta x_2^4 + x_1 + u \\
        &=: f(x) + u
\end{align*}
where we have defined the function $f: \R^2 \to \R$ to be unforced surface dynamics.

Letting $\hat{\theta} = 1$ be our best estimate of $\theta$, we can decompose
$f$ into an estimated component
\begin{equation*}
    \hat{f} = (a+1) x_1 + 3a x_2 + \hat{\theta} x_2^4
\end{equation*}
and an uncertain component $\tilde{f} = f - \hat{f}$ which is bounded by
\begin{equation}
    |\tilde{f}| \leq \max_{\theta} |\theta - \hat{\theta}| x_2^4 = 0.2 x_2^4
    \label{eq:error-bound}
\end{equation}
We select $u$ to eliminate the certain dynamics and dominate the uncertain
dynamics:
\begin{equation*}
    u = -\hat{f} - k \text{sgn}(s).
\end{equation*}

Then, taking the Lyapunov function $V(s) = 1/2 s^2$, we have
\begin{align}
    \dot{V}(s) = s \dot{s} = s \bigl[(\hat{f}(x) + \tilde{f}(x)) +
                                      (- \hat{f}(x) - k \text{sgn}(s))\bigr]
        &= \tilde{f}(x) s - k |s| \\
        &\leq |\tilde{f}(x) s| - k|s| \\
        &\leq (|\tilde{f}(x)| - k) |s|
        \label{eq:lyap-intermediate}
\end{align}
Choosing $k$ to be a function of $x_2$
\begin{equation*}
    k = k(x_2) = 0.2 x_2^4 + \eta
\end{equation*}
where $\eta > 0$ is some positive constant, by (\ref{eq:lyap-intermediate})
and (\ref{eq:error-bound}), we have
\begin{equation*}
    \dot{V} \leq - \eta |s| = - \eta \sqrt{2 \V(s)}
\end{equation*}

Therefore, the system converges to the surface in finite time, and
LaSalle's Invariance Principle shows that the origin is an asymptotically stable
equilibrium.











\end{document}
